%Select document class: 12 point, article
\documentclass[12pt]{article}

% This template is a combination of work done by Mike Catanzaro and Gabe Angelini-Knoll, both formerly of the WSU math dept, with some additions and synthesizations by Clayton Hayes (clayton.hayes@wayne.edu) for broader use
% Last updated 2018-03-10 by Aaron Willcock (ez9213@wayne.edu)

%AMS = American Mathematical Society
%math, symbols, theorems, fonts
\usepackage{amsmath, amssymb, amsthm, amsfonts}
\usepackage{times, flexisym, mdframed, xcolor}
\usepackage[left=1in,right=1in, top=1in, bottom=1in]{geometry}
\usepackage{ulem,multicol}


%%%%%%%%%%%%%%%%%%%%%%%%%%%%%%%%%%%%%%%%%%%%%%%%%%%%%%%%%%%%%%%%%%%%%%%%%%%%%%%%
%Commands definitions
\definecolor{backgroundcolorDef}{rgb}{0.1,0.1,0.12}
\definecolor{textcolorDef}{rgb}{0.77,0.77,0.77}
\definecolor{waynegreen}{RGB}{0,89,79}
\newcommand{\setbackgroundcolour}{\pagecolor{backgroundcolorDef}}  
\newcommand{\settextcolour}{\color{textcolorDef}}   
\newcommand{\invertbackgroundtext}{\setbackgroundcolour\settextcolour}
\mdfdefinestyle{darkAnswer}{%
    fontcolor=textcolorDef,
    backgroundcolor=backgroundcolorDef,
    linecolor=textcolorDef,
    % ----------------------------------
    frametitlebackgroundcolor=backgroundcolorDef,
    frametitlefontcolor=textcolorDef,
    frametitlerulecolor=waynegreen,
    frametitlerule=true, 
    frametitlerulewidth=2pt
    }
\mdfdefinestyle{darkQuesion}{%
    fontcolor=textcolorDef,
    backgroundcolor=backgroundcolorDef,
    linecolor=textcolorDef,
    linewidth=1pt
    }

%If this line is commented, then the appearance remains as usual.
\invertbackgroundtext
\newcommand{\lrb}[1]{\left[#1\right]}
\newcommand{\lrp}[1]{\left(#1\right)}
\newcommand{\lrs}[1]{\left\{#1\right\}}
\newcommand{\gof}[1]{g\lrp{#1}}
\newcommand{\fof}[1]{f\lrp{#1}}
\newcommand{\Zof}[1]{Z\lrp{#1}}
\newcommand{\gofinv}[1]{g^{-1}\lrp{#1}}
\newcommand{\fofinv}[1]{f^{-1}\lrp{#1}}
\newcommand{\Zofinv}[1]{Z^{-1}\lrp{#1}}
\newcommand{\nmod}[2]{#1\,\left(\text{mod}\,#2\right)}
\newcommand{\ngcd}[2]{\text{gcd}\left(#1\,,\,#2\right)}
\newcommand{\nord}[2]{\text{ord}_{#1}\left(#2\right)}
\newcommand{\dom}[1]{\text{dom}\left(#1\right)}
\newcommand{\ran}[1]{\text{ran}\left(#1\right)}
\newcommand{\degr}[1]{\text{deg}\left(#1\right)}
\newcommand{\N}{\mathbb{N}}
\newcommand{\Z}{\mathbb{Z}}
\newcommand{\Q}{\mathbb{Q}}
\newcommand{\R}{\mathbb{R}}
\newcommand{\C}{\mathbb{C}}
\newcommand{\abs}[1]{\left\vert#1\right\vert}
\newcommand{\Zm}[1]{\mathbb{Z}_{#1}}
\newcommand{\Zp}{\mathbb{Z}_p}
\newcommand{\numb}[1]{\noindent{\bf #1)}}
\newcommand{\inv}[1]{#1^{-1}}
\newcommand{\totient}[1]{\phi\left(#1\right)} 
\newcommand{\vhalfpg}{\vspace{5in}}
\newcommand{\vthirdpg}{\vspace{3in}}
\newcommand{\vquartpg}{\vspace{2in}}
\newcommand{\twocase}[2]{\begin{enumerate}
        \item[$ \implies$]{
            #1
        }
        \bigskip
        \item[$ \impliedby$]{
            #2
        }
    \end{enumerate}}
\newcommand{\msout}[1]{\text{\sout{\ensuremath{#1}}}}

%%%%%%%%%%%%%%%%%%%%%%%%%%%%%%%%%%%%%%%%%%%%%%%%%%%%%%%%%%%%%%%%%%%%%%%%%%%%%%%%
%Enhanced graphics support (https://ctan.org/pkg/graphicx?lang=en)
\usepackage{graphicx}

%Set default graphics path (replace 'figures/' with whatever directory your images are in)
%\graphicspath{{figures/}}


%Control layout of itemize, enumerate, description (https://ctan.org/pkg/enumitem?lang=en)
\usepackage{enumitem}

% Header and footer formatting options
\usepackage{fancyhdr}

%Control float placement. [section] = "stop floats at section boundaries is to change the definition of “\section” to include “\FloatBarrier”"
\usepackage[section]{placeins}

%Hypertext (links) in LaTeX. [option] = remove color and border on links.
\usepackage[hidelinks]{hyperref}

%\usepackage[all]{xy}
%\usepackage{mathtools}

%Pro­gram­ming fa­cil­i­ties for LaTeX class and pack­age authors
\usepackage{etoolbox}

%Indent first paragraph of every 'chapter' aka section
\usepackage{indentfirst}

%Title formatting option. [explicit] = make titles all caps
\usepackage[explicit]{titlesec}

%Standard package for selecting font encodings. [T1] = Support for accented characters
%(https://texfaq.org/FAQ-why-inp-font) - 
%(https://tex.stackexchange.com/questions/664/why-should-i-use-usepackaget1fontenc)
\usepackage[T1]{fontenc}

%Charter fonts
\usepackage{charter}

%\usepackage[expert]{mathdesign}

%Control table of contents, figures, etc
\usepackage{tocloft}

%Set space between lines. [option] = Double/single spacing necessary for properly formatting ToC/LoFT and titles.
\usepackage[doublespacing]{setspace}

%Create normal/logarithmic plots in two and three dimensions
\usepackage{pgfplots}
%Request a specific version
\pgfplotsset{compat=1.5}

\newcommand\scalemath[2]{\scalebox{#1}{\mbox{\ensuremath{\displaystyle #2}}}}

%Modify section headers (aka chapters) to be: centered, single space, large, bold, preceded with the word 'CHAPTER 'and with 0.4 em of space before title 
\titleformat{\section}[block]{\centering\singlespace\large\bf}{CHAPTER \thesection \hspace{0.4em} \MakeUppercase{#1}}{0em}{}{}

%Modify unnumbered section headers (aka chapters) to be: centered, single space, large, bold, preceded with the word 'CHAPTER 'and with 0.4 em of space before title 
%   We use unnumbered section headers for Acknowledgements, Dedication, Bio, etc. for consistency
\titleformat{name=\section,numberless}[block]{\centering\singlespace\large\bf}{\hspace{0.4em} \MakeUppercase{#1}}{0em}{}{}

%Remove spacing left of, above, and below sections and subsections
% \titlespacing{command}{left spacing}{before spacing}{after spacing}[right] (https://tex.stackexchange.com/questions/53338/reducing-spacing-after-headings)
\titlespacing{\section}{0pt}{0pt}{0pt}
\titlespacing{\subsection}{0pt}{0pt}{0pt}
\titlespacing{\subsubsection}{0pt}{0pt}{0pt}

% Modify TOC entries with tocloft: insert the word 'Chapter ' and then add 6em of space after the word
\renewcommand{\cftsecleader}{\cftdotfill{\cftdotsep}}
\renewcommand{\cftsecpresnum}{Chapter }
\cftsetindents{section}{0em}{6em}

% Call the other parts of the document 'part' and make the table of contents print them correctly
\renewcommand{\cftpartfont}{}
\renewcommand{\cftpartpagefont}{}
\renewcommand{\cftpartleader}{\cftdotfill{\cftdotsep}}
\renewcommand{\cftbeforepartskip}{0em}
\renewcommand{\cftpartindent}{1.5em}

% Change spacing above chapters in toc, remembering that they are really sections
\setlength{\cftbeforesecskip}{0em}

% Adds space after ToC entries to make them appear double-spaced
\renewcommand\cftsecafterpnum{\vskip\baselineskip}
\renewcommand\cftsubsecafterpnum{\vskip\baselineskip}
\renewcommand\cftsubsubsecafterpnum{\vskip\baselineskip}
\renewcommand\cftpartafterpnum{\vskip\baselineskip}

%Add space after LoF entries to make them appear double-spaced
\renewcommand\cftfigafterpnum{\vskip\baselineskip}

%Add space after LoT entries to make them appear double-spaced
\renewcommand\cfttabafterpnum{\vskip\baselineskip}

% Remove extra space above and below theorems, lemmas, props, etc.
%The important point in the following is: 0pt preskip and 0pt postskip
\makeatletter
\def\thm@space@setup{\thm@preskip=0pt
\thm@postskip=0pt}
\makeatother
\newtheoremstyle{newstyle}      
{} % Aboveskip 
{} % Below skip
{\mdseries} % Body font e.g.\mdseries,\bfseries,\scshape,\itshape
{} % Indent
{\bfseries}  % Head font e.g.\bfseries,\scshape,\itshape
{.} % Punctuation afer theorem header
{ } % Space after theorem header
{} % Heading

%The above does not fix the spacing around proof environments.
%Use the following to fix.
%The crucial point is "topsep0\p@", i.e., topsep = 0 pt.
%The rest is essentially copied from the standard AMS environment.
\makeatletter
\renewenvironment{proof}[1][\proofname]{\par
  \pushQED{\qed}%
  \normalfont \topsep0\p@\relax
  \trivlist
  \item[\hskip\labelsep\itshape
  #1\@addpunct{.}]\ignorespaces 
}{
  \popQED\endtrivlist\@endpefalse
}
\makeatother

%Remove default "References" text from \bibliography call
%   Removing this code will cause "References" to appear twice in the bibliography:
%       once for the \thebibliography call and once for the \section*{References} call
%   From: (https://tex.stackexchange.com/questions/132646/how-to-remove-the-references-title/132654)
\patchcmd{\thebibliography}{\section*{\refname}}{}{}{}

%Remove spacing between bibliography entries
%   From: (https://tex.stackexchange.com/questions/93859/condense-the-space-between-bibliographic-entries)
%   Set "OLDthebibliography" to be "thebibliography"
\let\OLDthebibliography\thebibliography
%   Renew command with 0pt spacing for paragraph and item separation
\renewcommand\thebibliography[1]{
    \OLDthebibliography{#1}
    \setlength{\parskip}{0pt}
    \setlength{\itemsep}{0pt plus 0.3ex}
}

%Formatting
%   Margins
\usepackage[left=1in,right=1in, top=1in, bottom=1in]{geometry}
%   \arraystretch = the FACTOR for spacing between two rows (Default = 1) 
\renewcommand{\arraystretch}{0.85}
%   \baselinestretch = The spacing between lines in a document (used for "double spacing")
%   From: (http://www.ctex.org/documents/latex/latex2e-html/baselineskip.html)
\renewcommand{\baselinestretch}{2}
%   \headrulewidth = "The thickness of a line under the header and above the footer"
%   From: (http://tug.ctan.org/macros/latex/contrib/fancyhdr/fancyhdr.pdf) pg 3
\renewcommand{\headrulewidth}{0pt}

%Define the StarTeX mathematical symbol
%From: (http://mirror.utexas.edu/ctan/macros/startex/startex.pdf)
\newcommand{\Mdef}[2]{\newcommand{#1}{\relax \ifmmode #2 \else $#2$\fi}}

%Spacing commands, which may or may not be used in the following.
\newcommand{\one}{{\rm \bf1}\hspace*{-0.035in} {\rm l}}
\newcommand{\nd}{\noindent}
\def\para#1{\vskip 0.4\baselineskip\noindent{\bf #1}}
\newcommand{\Vspc}{\vspace*{0.1in}}

\newcommand{\ID}{\index}
\makeatletter \@addtoreset{equation}{section} \makeatother

%Section Labeling
\renewcommand{\thesection}{\arabic{section}}

%Equation Numbering Ex. Section 3, Equation 2 = Equation 3.2
\renewcommand{\theequation}{\thesection.\arabic{equation}}

\renewcommand{\cftsecfont}{}
\renewcommand{\cftsecpagefont}{}

%List of Tables edits
\renewcommand{\cfttabpresnum}{Table }
\renewcommand{\cfttabindent}{1.0em}
\renewcommand{\cfttabnumwidth}{5.0em}

%List of Figures edits
\renewcommand{\cftfigpresnum}{Figure }
\renewcommand{\cftfigindent}{1.0em}
\renewcommand{\cftfignumwidth}{5.5em}

%Enter single spacing environment for toc, lot, and lof (see below)
\begin{document}
%Compile the title and jump to new page
%Create title (centered and bold font)
\centerline{\bf Algebra I}
\vspace{-0.4cm}
\centerline{\bf Winter 2020}

\begin{figure}[!htbp]
    \centering
    \includegraphics[width=0.25\linewidth]{fig/wsu_primary_stacked_color.pdf}
\end{figure}

\thispagestyle{empty}
\newpage
\setcounter{page}{2}
%Compile Chapter 1
\section{Integers}  
\subsection{Divisors}
%%%%%%%%%%%%%%%%%%%%%%%%%%%%%%%%%%%%%%%%%%%%%%%%%%%%%%%%%%%%%%%%%%%%%%%%%%%%%%%%
%%%%%%%%%%%%%%%%%%%%%%%%%%%%%%%%%%%%%%%%%%%%%%%%%%%%%%%%%%%%%%%%%%%%%%%%%%%%%%%%
%%%%%%%%%%%%%%%%%%%%%%%%%%%%%%%%%%%%%%%%%%%%%%%%%%%%%%%%%%%%%%%%%%%%%%%%%%%%%%%%
%%%%%%%%%%%%%%%%%%%%%%%%%%%%%%%%%%%%%%%%%%%%%%%%%%%%%%%%%%%%%%%%%%%%%%%%%%%%%%%%
\begin{mdframed}[style=darkQuesion]
1.    Let $m,n.r.s\in \Z$. If $m^2+n^2=r^2+s^2=mr+ns$, prove that $m=r$ and 
$n=s$.
\end{mdframed}

%%%%%%%%%%%%%%%%%%%%%%%%%%%%%%%%%%%%%%%%%%%%%%%%%%%%%%%%%%%%%%%%%%%%%%%%%%%%%%%%
\begin{mdframed}[style=darkAnswer,frametitle={Joe Starr}]
We select $m,n.r.s\in \Z$, given $m^2+n^2=r^2+s^2=mr+ns$ which can write
as $m^2+n^2-mr-ns=r^2+s^2-mr-ns$. From here we can simplify:
\begin{align*}
m^2+n^2-mr-ns=r^2+s^2-mr-ns &\Rightarrow 
m\lrp{m-r}+n\lrp{n-s}=r\lrp{r-m}+s\lrp{s-n} \\
&\Rightarrow m\lrp{m-r}+n\lrp{n-s}-r\lrp{r-m}-s\lrp{s-n}=0 \\
&\Rightarrow m\lrp{m-r}+r\lrp{m-r}+n\lrp{n-s}+s\lrp{n-s}=0 \\
&\Rightarrow \lrp{m-r}\lrp{m+r}+\lrp{n-s}\lrp{n+s}=0 \\
\end{align*}
from here we can see that in order for $\lrp{m-r}\lrp{m+r}+\lrp{n-s}\lrp{n+s}=0$
to be true $m=r$ and $n=s$.
\end{mdframed}
\newpage
%%%%%%%%%%%%%%%%%%%%%%%%%%%%%%%%%%%%%%%%%%%%%%%%%%%%%%%%%%%%%%%%%%%%%%%%%%%%%%%%
%%%%%%%%%%%%%%%%%%%%%%%%%%%%%%%%%%%%%%%%%%%%%%%%%%%%%%%%%%%%%%%%%%%%%%%%%%%%%%%%
%%%%%%%%%%%%%%%%%%%%%%%%%%%%%%%%%%%%%%%%%%%%%%%%%%%%%%%%%%%%%%%%%%%%%%%%%%%%%%%%
%%%%%%%%%%%%%%%%%%%%%%%%%%%%%%%%%%%%%%%%%%%%%%%%%%%%%%%%%%%%%%%%%%%%%%%%%%%%%%%%
\begin{mdframed}[style=darkQuesion]
3.    Find the quotient and reminder when $a$ id divided by $b$. 
\begin{itemize}
    \item [a] {$a=99$, $b=17$}
    \item [b] {$a=-99$, $b=17$}
    \item [c] {$a=17$, $b=99$}
    \item [d] {$a=-1017$, $b=99$}
\end{itemize}

\end{mdframed}

%%%%%%%%%%%%%%%%%%%%%%%%%%%%%%%%%%%%%%%%%%%%%%%%%%%%%%%%%%%%%%%%%%%%%%%%%%%%%%%%
\begin{mdframed}[style=darkAnswer,frametitle={Joe Starr}]
\begin{itemize}
    \item [a] {$99=17q+r\Rightarrow q=5, r=14$}
    \item [b] {$-99=17q+r\Rightarrow q=-6, r=3$}
    \item [c] {$17=99q+r\Rightarrow q=0, r=17$}
    \item [d] {$-1017=99q+r\Rightarrow q=-11, r=72$}
\end{itemize}
\end{mdframed}

\newpage
%%%%%%%%%%%%%%%%%%%%%%%%%%%%%%%%%%%%%%%%%%%%%%%%%%%%%%%%%%%%%%%%%%%%%%%%%%%%%%%%
%%%%%%%%%%%%%%%%%%%%%%%%%%%%%%%%%%%%%%%%%%%%%%%%%%%%%%%%%%%%%%%%%%%%%%%%%%%%%%%%
%%%%%%%%%%%%%%%%%%%%%%%%%%%%%%%%%%%%%%%%%%%%%%%%%%%%%%%%%%%%%%%%%%%%%%%%%%%%%%%%
%%%%%%%%%%%%%%%%%%%%%%%%%%%%%%%%%%%%%%%%%%%%%%%%%%%%%%%%%%%%%%%%%%%%%%%%%%%%%%%%
\begin{mdframed}[style=darkQuesion]
5.    Use the Euclidean algorithm to find the following greatest common divisors
\begin{multicols}{2}
    \begin{itemize}
        \item [a] {$\lrp{6643,2873}$}
        \item [b] {$\lrp{7684,4148}$}
        \item [c] {$\lrp{26460,12600}$}
        \item [d] {$\lrp{6540,1206}$}
        \item [e] {$\lrp{12091,8439}$}
    \end{itemize}
\end{multicols}
\end{mdframed}

%%%%%%%%%%%%%%%%%%%%%%%%%%%%%%%%%%%%%%%%%%%%%%%%%%%%%%%%%%%%%%%%%%%%%%%%%%%%%%%%
\begin{mdframed}[style=darkAnswer,frametitle={Joe Starr}]
\begin{multicols}{2}
\begin{itemize}
    \item [(a)] {$\lrp{6643,2873}$
    \begin{align*}
        6643&=2873*2+897\\
        2873&=897*3+182\\
        897&=182*4+169\\
        182&=169*1+13\\
        169&=13*13\\
    \end{align*}
    }
    \item [(b)] {$\lrp{7684,4148}$
    \begin{align*}
        7684&=4148*1+3536\\
        4148&=3536*1+612\\
        3536&=612*5+476\\
        612&=476*1+136\\
        476&=136*3+68\\
        136&=68*68\\
    \end{align*}
    }
    \item [(c)] {$\lrp{26460,12600}$
    \begin{align*}
        26460&=12600*2+1260\\
        12600&=1260*10\\
    \end{align*}
    }
\end{itemize}
\end{multicols}
    \begin{multicols}{2}
        \begin{itemize}
    \item [(d)] {$\lrp{6540,1206}$
    \begin{align*}
        6540&=1206*5+510\\
        1206&=510*2+186\\
        510&=186*2+138\\
        186&=138*1+48\\
        138&=48*2+42\\
        48&=42*1+6\\
        42&=6*7\\
    \end{align*}
    }
    \item [(e)] {$\lrp{12091,8439}$
    \begin{align*}
        12091&=8439*1+3652\\
        8439&=3652*2+1135\\
        3652&=1135*3+247\\
        1135&=247*4+147\\
        247&=147*1+100\\
        147&=100*1+47\\
        100&=47*2+6\\
        47&=6*7+5\\
        6&=5*1+1\\
        5&=1*5\\
    \end{align*}
    }
\end{itemize}
\end{multicols}
\end{mdframed}
\newpage
%%%%%%%%%%%%%%%%%%%%%%%%%%%%%%%%%%%%%%%%%%%%%%%%%%%%%%%%%%%%%%%%%%%%%%%%%%%%%%%%
%%%%%%%%%%%%%%%%%%%%%%%%%%%%%%%%%%%%%%%%%%%%%%%%%%%%%%%%%%%%%%%%%%%%%%%%%%%%%%%%
%%%%%%%%%%%%%%%%%%%%%%%%%%%%%%%%%%%%%%%%%%%%%%%%%%%%%%%%%%%%%%%%%%%%%%%%%%%%%%%%
%%%%%%%%%%%%%%%%%%%%%%%%%%%%%%%%%%%%%%%%%%%%%%%%%%%%%%%%%%%%%%%%%%%%%%%%%%%%%%%%
\begin{mdframed}[style=darkQuesion]
    7.    For each part of Exercise 5, find integers $m$ and $n$ such that 
$\lrp{a,b}$ is expressed in the form $ma+nb$. 

\end{mdframed}

%%%%%%%%%%%%%%%%%%%%%%%%%%%%%%%%%%%%%%%%%%%%%%%%%%%%%%%%%%%%%%%%%%%%%%%%%%%%%%%%
\begin{mdframed}[style=darkAnswer,frametitle={Joe Starr}]
\begin{itemize}
    \item [(a)] {$\lrp{6643,2873}$\\
        $\lrp{6643}-16+\lrp{2873}37=13$
    }
    \item [(b)] {$\lrp{7684,4148}$\\
    $\lrp{7684}27+\lrp{4148}-50=68$
    }
    \item [(c)] {$\lrp{26460,12600}$
    $\lrp{26460}1+\lrp{12600}-2=1260$
    }
    \item [(d)] {$\lrp{6540,1206}$
   $\lrp{6540}-26+\lrp{1206}141=6$
    }
    \item [(e)] {$\lrp{12091,8439}$
    $\lrp{12091}1435+\lrp{8439}-2056=1$
    }
\end{itemize}
\end{mdframed}
\newpage
%%%%%%%%%%%%%%%%%%%%%%%%%%%%%%%%%%%%%%%%%%%%%%%%%%%%%%%%%%%%%%%%%%%%%%%%%%%%%%%%
%%%%%%%%%%%%%%%%%%%%%%%%%%%%%%%%%%%%%%%%%%%%%%%%%%%%%%%%%%%%%%%%%%%%%%%%%%%%%%%%
%%%%%%%%%%%%%%%%%%%%%%%%%%%%%%%%%%%%%%%%%%%%%%%%%%%%%%%%%%%%%%%%%%%%%%%%%%%%%%%%
%%%%%%%%%%%%%%%%%%%%%%%%%%%%%%%%%%%%%%%%%%%%%%%%%%%%%%%%%%%%%%%%%%%%%%%%%%%%%%%%
\begin{mdframed}[style=darkQuesion]
9.  let $a,b,c$ be integers such that $a+b+c=0$. Show that if $n$ is an integer
which is a divisor of two of the three integers, then it is also a divisor of 
the third. 

\end{mdframed}

%%%%%%%%%%%%%%%%%%%%%%%%%%%%%%%%%%%%%%%%%%%%%%%%%%%%%%%%%%%%%%%%%%%%%%%%%%%%%%%%
\begin{mdframed}[style=darkAnswer,frametitle={Joe Starr}]
Select $a,b,c\in \Z$ to satisfy $a+b+c=0$, WLOG let $n\in \Z$ such that 
$n\vert a$ and $n\vert b$. Since $\lrp{a+b}+c=0$ it must be that $\lrp{a+b}=-c$.
From here we must show $n\vert\lrp{a+b}$, or $a+b=nq$. Since $n\vert a$ and 
$n\vert b$ we may write $a=nq_1$ and $b=nq_2$, yielding, 
$nq_1+nq_2=n\lrp{q_1+q_2}=nq$ thus $n\vert c$, as desired.$\square$ 
\end{mdframed}
\newpage
%%%%%%%%%%%%%%%%%%%%%%%%%%%%%%%%%%%%%%%%%%%%%%%%%%%%%%%%%%%%%%%%%%%%%%%%%%%%%%%%
%%%%%%%%%%%%%%%%%%%%%%%%%%%%%%%%%%%%%%%%%%%%%%%%%%%%%%%%%%%%%%%%%%%%%%%%%%%%%%%%
%%%%%%%%%%%%%%%%%%%%%%%%%%%%%%%%%%%%%%%%%%%%%%%%%%%%%%%%%%%%%%%%%%%%%%%%%%%%%%%%
%%%%%%%%%%%%%%%%%%%%%%%%%%%%%%%%%%%%%%%%%%%%%%%%%%%%%%%%%%%%%%%%%%%%%%%%%%%%%%%%
\begin{mdframed}[style=darkQuesion]
13.  Show that if $n$ is any integer, then $\lrp{10n_3,5n+2}=1$
\end{mdframed}

%%%%%%%%%%%%%%%%%%%%%%%%%%%%%%%%%%%%%%%%%%%%%%%%%%%%%%%%%%%%%%%%%%%%%%%%%%%%%%%%
\begin{mdframed}[style=darkAnswer,frametitle={Joe Starr}]
We begin with the Euclidean algorithm, 
\begin{align*}
    10n+3&=\lrp{5n+2}1+\lrp{5n+1}\\
    5n+2&=\lrp{5n+1}1+1\\
\end{align*}
from here we have $\lrp{10n+3,5n+2}=\lrp{5n+2,5n+1}=1$, as desired.
\end{mdframed}
\newpage
%%%%%%%%%%%%%%%%%%%%%%%%%%%%%%%%%%%%%%%%%%%%%%%%%%%%%%%%%%%%%%%%%%%%%%%%%%%%%%%%
%%%%%%%%%%%%%%%%%%%%%%%%%%%%%%%%%%%%%%%%%%%%%%%%%%%%%%%%%%%%%%%%%%%%%%%%%%%%%%%%
%%%%%%%%%%%%%%%%%%%%%%%%%%%%%%%%%%%%%%%%%%%%%%%%%%%%%%%%%%%%%%%%%%%%%%%%%%%%%%%%
%%%%%%%%%%%%%%%%%%%%%%%%%%%%%%%%%%%%%%%%%%%%%%%%%%%%%%%%%%%%%%%%%%%%%%%%%%%%%%%%
\begin{mdframed}[style=darkQuesion]
15.  For what positive integers $n$ is it true that $\lrp{n,n+2}=2$? Prove your
claim.
\end{mdframed}

%%%%%%%%%%%%%%%%%%%%%%%%%%%%%%%%%%%%%%%%%%%%%%%%%%%%%%%%%%%%%%%%%%%%%%%%%%%%%%%%
\begin{mdframed}[style=darkAnswer,frametitle={Joe Starr}]
The conjecture is that the statement is true for even values of $n$.
We begin with rewriting $n$ in terms of $k$, $n=2k$the Euclidean algorithm, 
\begin{align*}
    \lrp{2k}+2&=\lrp{2k}1+\lrp{2}\\
    2k&=\lrp{2}{k}\\
\end{align*}
from here we have $\lrp{n+2,n}=\lrp{2k+2,2k}=2$, as desired.
\end{mdframed}
\newpage
%%%%%%%%%%%%%%%%%%%%%%%%%%%%%%%%%%%%%%%%%%%%%%%%%%%%%%%%%%%%%%%%%%%%%%%%%%%%%%%%
%%%%%%%%%%%%%%%%%%%%%%%%%%%%%%%%%%%%%%%%%%%%%%%%%%%%%%%%%%%%%%%%%%%%%%%%%%%%%%%%
%%%%%%%%%%%%%%%%%%%%%%%%%%%%%%%%%%%%%%%%%%%%%%%%%%%%%%%%%%%%%%%%%%%%%%%%%%%%%%%%
%%%%%%%%%%%%%%%%%%%%%%%%%%%%%%%%%%%%%%%%%%%%%%%%%%%%%%%%%%%%%%%%%%%%%%%%%%%%%%%%
\begin{mdframed}[style=darkQuesion]
17.  Show that the positive integer $k$ is the difference of two odd squares if 
and only if $k$ is divisible by $8$.
\end{mdframed}

%%%%%%%%%%%%%%%%%%%%%%%%%%%%%%%%%%%%%%%%%%%%%%%%%%%%%%%%%%%%%%%%%%%%%%%%%%%%%%%%
\begin{mdframed}[style=darkAnswer,frametitle={Joe Starr}]

We begin by writing $k=a^2-b^2$, since $a$ and $b$ are odd we can write, 
\begin{align*}
    a&=2r+1\\
    b&=2s+1
\end{align*} 
from here we have $q^2-b^2=4\lrp{r+s+1}\lrp{r-s}$. Since $k>0$ we must consider 
two cases $r-s=2m+1$ and $r-s=2m$. 
\begin{itemize}[align=left]
    \item [$r-s=2m$:]{\hspace{.5in}\newline
        In this case we have $q^2-b^2=4\lrp{r+s+1}2m=8\lrp{r+s+1}m$ and we are 
        done. 
    }
    \item [$r-s=2m+1$:]{\hspace{.5in}\newline
        In this case we have $r-s=2m+1$ and $r+s=r-s+2s=2m+1+2s$
        \begin{align*}
            q^2-b^2&=4\lrp{r+s+1}\lrp{2m+1}\\
            &=4\lrp{2m\lrp{r+s+1}+\lrp{r+s+1}}\\            
            &=4\lrp{\lrp{2mr+2ms+2m}+\lrp{r+s+1}}\\            
            &=4\lrp{2mr+2ms+2m+r+s+1}\\         
            &=4\lrp{2mr+2ms+2m+2m+1+2s+1}\\         
            &=4\lrp{2mr+2ms+2m+2m+2s+2}\\         
            &=8\lrp{mr+ms+m+m+s+1}\\         
        \end{align*}
        as desired.
    }
\end{itemize}

\end{mdframed}
\newpage
%%%%%%%%%%%%%%%%%%%%%%%%%%%%%%%%%%%%%%%%%%%%%%%%%%%%%%%%%%%%%%%%%%%%%%%%%%%%%%%%
%%%%%%%%%%%%%%%%%%%%%%%%%%%%%%%%%%%%%%%%%%%%%%%%%%%%%%%%%%%%%%%%%%%%%%%%%%%%%%%%
%%%%%%%%%%%%%%%%%%%%%%%%%%%%%%%%%%%%%%%%%%%%%%%%%%%%%%%%%%%%%%%%%%%%%%%%%%%%%%%%
%%%%%%%%%%%%%%%%%%%%%%%%%%%%%%%%%%%%%%%%%%%%%%%%%%%%%%%%%%%%%%%%%%%%%%%%%%%%%%%%
\subsection{Primes}
%%%%%%%%%%%%%%%%%%%%%%%%%%%%%%%%%%%%%%%%%%%%%%%%%%%%%%%%%%%%%%%%%%%%%%%%%%%%%%%%
%%%%%%%%%%%%%%%%%%%%%%%%%%%%%%%%%%%%%%%%%%%%%%%%%%%%%%%%%%%%%%%%%%%%%%%%%%%%%%%%
%%%%%%%%%%%%%%%%%%%%%%%%%%%%%%%%%%%%%%%%%%%%%%%%%%%%%%%%%%%%%%%%%%%%%%%%%%%%%%%%
%%%%%%%%%%%%%%%%%%%%%%%%%%%%%%%%%%%%%%%%%%%%%%%%%%%%%%%%%%%%%%%%%%%%%%%%%%%%%%%%
\begin{mdframed}[style=darkQuesion]
1. Find the prime factorizations of each of the following numbers, and use the 
them to compute the greatest common divisor and least common multiple of the 
given pairs of numbers. 
\begin{multicols}{3}
\begin{itemize}
    \item [(a)] {$\lrp{35,14}$
    }
    \item [(b)] {$\lrp{15,11}$
    }
    \item [(c)] {$\lrp{252,11}$
    }
    \item [(d)] {$\lrp{7684,4148}$
    }
    \item [(e)] {$\lrp{6643,2873}$
    }
\end{itemize}
\end{multicols}
\end{mdframed}

%%%%%%%%%%%%%%%%%%%%%%%%%%%%%%%%%%%%%%%%%%%%%%%%%%%%%%%%%%%%%%%%%%%%%%%%%%%%%%%%
\begin{mdframed}[style=darkAnswer,frametitle={Joe Starr}]
    \begin{multicols}{2}
    \begin{itemize}
        \item [(a)] {
        \begin{multicols}{2}
        $\lrp{35,14}$\\
        $35: 5,7$ \\
        $14: 2,7$ \\
        gcd: $7$ \\
        lcm: $70$ 
        \end{multicols}
        }
        \item [(b)] {
            \begin{multicols}{2}
        $\lrp{15,11}$ \\
        $15: 3,5$ \\
        $11: 11$ \\
        gcd: $1$ \\
        lcm: $165$ 
    \end{multicols}
        }
        \item [(c)] {
            \begin{multicols}{2}
        $\lrp{252,180}$ \\
        $252: 2,2,3,3,7$ \\
        $180: 2,2,3,3,5$ \\
        gcd: $36$ \\
        lcm: $1260$ 
    \end{multicols}
        }
        \item [(d)] {
            \begin{multicols}{2}
        $\lrp{7684,4148}$ \\
        $7684: 2,2,17,113$ \\
        $4148: 2,2,17,61$ \\
        gcd: $68$ \\
        lcm: $468724$ 
    \end{multicols}
        }
        \item [(e)] {
            \begin{multicols}{2}
        $\lrp{6643,2873}$ \\
        $6643: 7,13,73$ \\
        $2873: 13,13,17$ \\
        gcd: $13$ \\
        lcm: $1468103$ 
    \end{multicols}
        }
    \end{itemize}
\end{multicols}
\end{mdframed}
\newpage
%%%%%%%%%%%%%%%%%%%%%%%%%%%%%%%%%%%%%%%%%%%%%%%%%%%%%%%%%%%%%%%%%%%%%%%%%%%%%%%%
%%%%%%%%%%%%%%%%%%%%%%%%%%%%%%%%%%%%%%%%%%%%%%%%%%%%%%%%%%%%%%%%%%%%%%%%%%%%%%%%
%%%%%%%%%%%%%%%%%%%%%%%%%%%%%%%%%%%%%%%%%%%%%%%%%%%%%%%%%%%%%%%%%%%%%%%%%%%%%%%%
%%%%%%%%%%%%%%%%%%%%%%%%%%%%%%%%%%%%%%%%%%%%%%%%%%%%%%%%%%%%%%%%%%%%%%%%%%%%%%%%
\begin{mdframed}[style=darkQuesion]
2. US the sieve of Eratosthenes to find all prime numbers less than 200.
\end{mdframed}

%%%%%%%%%%%%%%%%%%%%%%%%%%%%%%%%%%%%%%%%%%%%%%%%%%%%%%%%%%%%%%%%%%%%%%%%%%%%%%%%
\begin{mdframed}[style=darkAnswer,frametitle={Joe Starr}]
    \begin{center}
        \begin{tabular}{| c | c | c | c | c | c | c | c | c | c |}
        \hline
        $ $ & $2$ & $3$ & $\msout{4}$ & $5$ & $\msout{6}$ & $7$ & $\msout{8}$ & 
        $\msout{9}$ & $\msout{10}$ \\
        \hline
        $11$ & $\msout{12}$ & $13$ & $\msout{14}$ & $\msout{15}$ & $\msout{16}$ 
        & $17$ & $\msout{18}$ & $19$ & $\msout{20}$ \\
        \hline
        $\msout{21}$ & $\msout{22}$ & $23$ & $\msout{24}$ & $\msout{25}$ & 
        $\msout{26}$ & $\msout{27}$ & $\msout{28}$ & $29$ & $\msout{30}$ \\
        \hline
        $31$ & $\msout{32}$ & $\msout{33}$ & $\msout{34}$ & $\msout{35}$ & 
        $\msout{36}$ & $37$ & $\msout{38}$ & $\msout{39}$ & $\msout{40}$ \\
        \hline
        $41$ & $\msout{42}$ & $43$ & $\msout{44}$ & $\msout{45}$ & $\msout{46}$ 
        & $47$ & $\msout{48}$ & $\msout{49}$ & $\msout{50}$ \\
        \hline
        $\msout{51}$ & $\msout{52}$ & $53$ & $\msout{54}$ & $\msout{55}$ & 
        $\msout{56}$ & $\msout{57}$ & $\msout{58}$ & $59$ & $\msout{60}$ \\
        \hline
        $61$ & $\msout{62}$ & $\msout{63}$ & $\msout{64}$ & $\msout{65}$ & 
        $\msout{66}$ & $67$ & $\msout{68}$ & $\msout{69}$ & $\msout{70}$ \\
        \hline
        $71$ & $\msout{72}$ & $73$ & $\msout{74}$ & $\msout{75}$ & $\msout{76}$ 
        & $\msout{77}$ & $\msout{78}$ & $79$ & $\msout{80}$ \\
        \hline
        $\msout{81}$ & $\msout{82}$ & $83$ & $\msout{84}$ & $\msout{85}$ & 
        $\msout{86}$ & $\msout{87}$ & $\msout{88}$ & $89$ & $\msout{90}$ \\
        \hline
        $\msout{91}$ & $\msout{92}$ & $\msout{93}$ & $\msout{94}$ & $\msout{95}$
        & $\msout{96}$ & $97$ & $\msout{98}$ & $\msout{99}$ & $\msout{100}$ \\
        \hline
        $101$ & $\msout{102}$ & $103$ & $\msout{104}$ & $\msout{105}$ & 
        $\msout{106}$ & $107$ & $\msout{108}$ & $109$ & $\msout{110}$ \\
        \hline
        $\msout{111}$ & $\msout{112}$ & $113$ & $\msout{114}$ & $\msout{115}$ & 
        $\msout{116}$ & $\msout{117}$ & $\msout{118}$ & $\msout{119}$ & 
        $\msout{120}$ \\
        \hline
        $\msout{121}$ & $\msout{122}$ & $\msout{123}$ & $\msout{124}$ & 
        $\msout{125}$ & $\msout{126}$ & $127$ & $\msout{128}$ & $\msout{129}$ & 
        $\msout{130}$ \\
        \hline
        $131$ & $\msout{132}$ & $\msout{133}$ & $\msout{134}$ & $\msout{135}$ & 
        $\msout{136}$ & $137$ & $\msout{138}$ & $139$ & $\msout{140}$ \\
        \hline
        $\msout{141}$ & $\msout{142}$ & $\msout{143}$ & $\msout{144}$ & 
        $\msout{145}$ & $\msout{146}$ & $\msout{147}$ & $\msout{148}$ & $149$ & 
        $\msout{150}$ \\
        \hline
        $151$ & $\msout{152}$ & $\msout{153}$ & $\msout{154}$ & $\msout{155}$ & 
        $\msout{156}$ & $157$ & $\msout{158}$ & $\msout{159}$ & $\msout{160}$ \\
        \hline
        $\msout{161}$ & $\msout{162}$ & $163$ & $\msout{164}$ & $\msout{165}$ & 
        $\msout{166}$ & $167$ & $\msout{168}$ & $\msout{169}$ & $\msout{170}$ \\
        \hline
        $\msout{171}$ & $\msout{172}$ & $173$ & $\msout{174}$ & $\msout{175}$ & 
        $\msout{176}$ & $\msout{177}$ & $\msout{178}$ & $179$ & $\msout{180}$ \\
        \hline
        $181$ & $\msout{182}$ & $\msout{183}$ & $\msout{184}$ & $\msout{185}$ & 
        $\msout{186}$ & $\msout{187}$ & $\msout{188}$ & $\msout{189}$ & 
        $\msout{190}$ \\
        \hline
        $191$ & $\msout{192}$ & $193$ & $\msout{194}$ & $\msout{195}$ & 
        $\msout{196}$ & $197$ & $\msout{198}$ & $199$ & $\msout{200}$\\
        \hline
        \end{tabular}
        \end{center}    
\end{mdframed}
\newpage
%%%%%%%%%%%%%%%%%%%%%%%%%%%%%%%%%%%%%%%%%%%%%%%%%%%%%%%%%%%%%%%%%%%%%%%%%%%%%%%%
%%%%%%%%%%%%%%%%%%%%%%%%%%%%%%%%%%%%%%%%%%%%%%%%%%%%%%%%%%%%%%%%%%%%%%%%%%%%%%%%
%%%%%%%%%%%%%%%%%%%%%%%%%%%%%%%%%%%%%%%%%%%%%%%%%%%%%%%%%%%%%%%%%%%%%%%%%%%%%%%%
%%%%%%%%%%%%%%%%%%%%%%%%%%%%%%%%%%%%%%%%%%%%%%%%%%%%%%%%%%%%%%%%%%%%%%%%%%%%%%%%
\begin{mdframed}[style=darkQuesion]
3. For each composite number $a$. with $4\leq a\leq 20$, find all positive 
numbers less than $a$ that are relatively prime to $a$.
\end{mdframed}

%%%%%%%%%%%%%%%%%%%%%%%%%%%%%%%%%%%%%%%%%%%%%%%%%%%%%%%%%%%%%%%%%%%%%%%%%%%%%%%%
\begin{mdframed}[style=darkAnswer,frametitle={Joe Starr}]
\begin{multicols}{1}
\begin{itemize}
    \item[$4:$]  {$2, 3$}
    \item[$6:$]  {$2, 3, 5$}
    \item[$8:$]  {$2, 3, 5, 7$}
    \item[$9:$]  {$2, 3, 4, 5, 7, 8$}
    \item[$10:$] {$2, 3, 5, 7, 9$}
    \item[$12:$] {$2, 3, 5, 7, 11$}
    \item[$14:$] {$2, 3, 5, 7, 9, 11, 13$}
    \item[$15:$] {$2, 3, 4, 5, 7, 8, 11, 13, 14$}
    \item[$16:$] {$2, 3, 5, 7, 9, 11, 13, 15$}
    \item[$18:$] {$2, 3, 5, 7, 11, 13, 17$}
    \item[$20:$] {$2, 3, 5, 7, 9, 11, 13, 17, 19$}
\end{itemize}
\end{multicols}
\end{mdframed}
\newpage
%%%%%%%%%%%%%%%%%%%%%%%%%%%%%%%%%%%%%%%%%%%%%%%%%%%%%%%%%%%%%%%%%%%%%%%%%%%%%%%%
%%%%%%%%%%%%%%%%%%%%%%%%%%%%%%%%%%%%%%%%%%%%%%%%%%%%%%%%%%%%%%%%%%%%%%%%%%%%%%%%
%%%%%%%%%%%%%%%%%%%%%%%%%%%%%%%%%%%%%%%%%%%%%%%%%%%%%%%%%%%%%%%%%%%%%%%%%%%%%%%%
%%%%%%%%%%%%%%%%%%%%%%%%%%%%%%%%%%%%%%%%%%%%%%%%%%%%%%%%%%%%%%%%%%%%%%%%%%%%%%%%
\begin{mdframed}[style=darkQuesion]
4. Find all positive integers less than 60 and relatively prime to $60$.
\end{mdframed}

%%%%%%%%%%%%%%%%%%%%%%%%%%%%%%%%%%%%%%%%%%%%%%%%%%%%%%%%%%%%%%%%%%%%%%%%%%%%%%%%
\begin{mdframed}[style=darkAnswer,frametitle={Joe Starr}]
\begin{itemize}
    \item[$60:$] {$2, 3, 5, 7, 11, 13, 17, 19, 23, 29, 31, 
                    37, 41, 43, 47, 49, 53, 59 $}
\end{itemize}
\end{mdframed}
\newpage
%%%%%%%%%%%%%%%%%%%%%%%%%%%%%%%%%%%%%%%%%%%%%%%%%%%%%%%%%%%%%%%%%%%%%%%%%%%%%%%%
%%%%%%%%%%%%%%%%%%%%%%%%%%%%%%%%%%%%%%%%%%%%%%%%%%%%%%%%%%%%%%%%%%%%%%%%%%%%%%%%
%%%%%%%%%%%%%%%%%%%%%%%%%%%%%%%%%%%%%%%%%%%%%%%%%%%%%%%%%%%%%%%%%%%%%%%%%%%%%%%%
%%%%%%%%%%%%%%%%%%%%%%%%%%%%%%%%%%%%%%%%%%%%%%%%%%%%%%%%%%%%%%%%%%%%%%%%%%%%%%%%
\begin{mdframed}[style=darkQuesion]
\begin{itemize}
    \item [9. (a)] {For which $n\in \Z^{\text{+}}$ is $n^{3}-1$ a prime number?}
    \item [(b)] {For which $n\in \Z^{\text{+}}$ is $n^{3}+1$ a prime number?}
    \item [(c)] {For which $n\in \Z^{\text{+}}$ is $n^{2}-1$ a prime number?}
    \item [(d)] {For which $n\in \Z^{\text{+}}$ is $n^{2}+1$ a prime number?}
    \end{itemize}
\end{mdframed}

%%%%%%%%%%%%%%%%%%%%%%%%%%%%%%%%%%%%%%%%%%%%%%%%%%%%%%%%%%%%%%%%%%%%%%%%%%%%%%%%
\begin{mdframed}[style=darkAnswer,frametitle={Joe Starr}]
    \begin{itemize}
        \item [(a)] {We can factor $n^{3}-1$ into $(n - 1) (n^2 + n + 1)$. We 
        have then $n-1\vert n^{3}-1$, for $n^{3}-1$ to be prime $n-1$ must be 
        $1$. This happens only for $n=2$.
        }
        \item [(b)] {We can factor $n^{3}+1$ into $(n + 1) (n^2 - n + 1)$. We 
        have then $(n^2 - n + 1)\vert n^{3}+1$, 
        for $n^{3}+1$ to be prime $(n^2 - n + 1)$ must be $1$. 
        This happens only for $n=1$. }
        \item [(c)] {We can factor $n^{2}-1$ into $(n - 1) (n + 1)$. We 
        have then $(n 1 1)\vert n^{2}-1$, 
        for $n^{2}-1$ to be prime $(n - 1)$ must be $1$. 
        This happens only for $n=2$.
        For which $n\in \Z^{\text{+}}$ is $n^{2}-1$ a prime number?}
        \item [(d)] {????
        }
    \end{itemize}
\end{mdframed}

\newpage
%%%%%%%%%%%%%%%%%%%%%%%%%%%%%%%%%%%%%%%%%%%%%%%%%%%%%%%%%%%%%%%%%%%%%%%%%%%%%%%%
%%%%%%%%%%%%%%%%%%%%%%%%%%%%%%%%%%%%%%%%%%%%%%%%%%%%%%%%%%%%%%%%%%%%%%%%%%%%%%%%
%%%%%%%%%%%%%%%%%%%%%%%%%%%%%%%%%%%%%%%%%%%%%%%%%%%%%%%%%%%%%%%%%%%%%%%%%%%%%%%%
%%%%%%%%%%%%%%%%%%%%%%%%%%%%%%%%%%%%%%%%%%%%%%%%%%%%%%%%%%%%%%%%%%%%%%%%%%%%%%%%
\begin{mdframed}[style=darkQuesion]
11. Prove that $n^4+4^n$ is composite if $n>1$.
\end{mdframed}

%%%%%%%%%%%%%%%%%%%%%%%%%%%%%%%%%%%%%%%%%%%%%%%%%%%%%%%%%%%%%%%%%%%%%%%%%%%%%%%%
\begin{mdframed}[style=darkAnswer,frametitle={Joe Starr}]
We are presented with two potability's, $n$ is even or $n$ is odd.
\begin{itemize}[align=left]
    \item [$n$ even]{\hspace{.5in}\newline
        It's obvious that $n^4+4^n$ is an even not $2$ and can't be prime. 
    }
    \item [$n$ odd]{\hspace{.5in}\newline
    We begin by completing the square 
    \begin{align*}
        n^4+4^n&= n^4+4^n\\
        &=\lrp{n^2}^2+\lrp{2^n}^2\\
        &=\lrp{n^2+2^n}^2-2n^2 2^n
    \end{align*}
    We from here we observe that $2^n2=2^{n+1}$, since $n$ is odd $n+1$ is even 
    yielding $2^{n+1}=2^{2k}$. We can see we have a difference of squares
    \begin{align*}
        \lrp{n^2+2^n}^2-2n^2 2^n&=\lrp{n^2+2^n}^2-\lrp{2^nn}^2\\
        &=\lrp{n^2+2^n+2^nn}\lrp{n^2+2^n-2^nn}
    \end{align*}
    since we are restricted to $n>1$ we can see that both $\lrp{n^2+2^n+2^nn}>1$
    and $\lrp{n^2+2^n-2^nn}>1$ for all $n$. Making $n^4+4^n$ composite as 
    desired. 
    }
\end{itemize}
\end{mdframed}
\newpage
%%%%%%%%%%%%%%%%%%%%%%%%%%%%%%%%%%%%%%%%%%%%%%%%%%%%%%%%%%%%%%%%%%%%%%%%%%%%%%%%
%%%%%%%%%%%%%%%%%%%%%%%%%%%%%%%%%%%%%%%%%%%%%%%%%%%%%%%%%%%%%%%%%%%%%%%%%%%%%%%%
%%%%%%%%%%%%%%%%%%%%%%%%%%%%%%%%%%%%%%%%%%%%%%%%%%%%%%%%%%%%%%%%%%%%%%%%%%%%%%%%
%%%%%%%%%%%%%%%%%%%%%%%%%%%%%%%%%%%%%%%%%%%%%%%%%%%%%%%%%%%%%%%%%%%%%%%%%%%%%%%%
\begin{mdframed}[style=darkQuesion]
13. Let $a,b,c$  be positive integers, and let $d=\lrp{a,b}$. Since $d\vert a$, 
there exists an integer $h$ with $a=dh$. Show that $a\vert bc$, then $h\vert c$. 
\end{mdframed}

%%%%%%%%%%%%%%%%%%%%%%%%%%%%%%%%%%%%%%%%%%%%%%%%%%%%%%%%%%%%%%%%%%%%%%%%%%%%%%%%
\begin{mdframed}[style=darkAnswer,frametitle={Joe Starr}]
    
\end{mdframed}
\newpage
%%%%%%%%%%%%%%%%%%%%%%%%%%%%%%%%%%%%%%%%%%%%%%%%%%%%%%%%%%%%%%%%%%%%%%%%%%%%%%%%
%%%%%%%%%%%%%%%%%%%%%%%%%%%%%%%%%%%%%%%%%%%%%%%%%%%%%%%%%%%%%%%%%%%%%%%%%%%%%%%%
%%%%%%%%%%%%%%%%%%%%%%%%%%%%%%%%%%%%%%%%%%%%%%%%%%%%%%%%%%%%%%%%%%%%%%%%%%%%%%%%
%%%%%%%%%%%%%%%%%%%%%%%%%%%%%%%%%%%%%%%%%%%%%%%%%%%%%%%%%%%%%%%%%%%%%%%%%%%%%%%%
\begin{mdframed}[style=darkQuesion]
14. Show that $a\Z \cap b\Z=\lrb{a,b}$.
\end{mdframed}

%%%%%%%%%%%%%%%%%%%%%%%%%%%%%%%%%%%%%%%%%%%%%%%%%%%%%%%%%%%%%%%%%%%%%%%%%%%%%%%%
\begin{mdframed}[style=darkAnswer,frametitle={Joe Starr}]
    Let $x\in \lrp{a\Z \cap b\Z}$, since $x\in a\Z$ we have $x=aq_1$, similarly 
    since $x\in b\Z$ we have$x=bq_2$. We can see that $x=abq$, this means $x$ 
    is a multiple of $\lrb{a,b}$ putting $x\in \lrb{a,b}$. Next, we let 
    $x\in \lrb{a,b}\Z$, this means $x$ is of the form $x=\lrb{a,b}q$. We can see
    that $a\vert x$ and $b\vert x$ since $a\vert \lrb{a,b}$, This makes 
    $x\in a\Z$ and $x\in b\Z$, as desired. 
\end{mdframed}
\newpage
%%%%%%%%%%%%%%%%%%%%%%%%%%%%%%%%%%%%%%%%%%%%%%%%%%%%%%%%%%%%%%%%%%%%%%%%%%%%%%%%
%%%%%%%%%%%%%%%%%%%%%%%%%%%%%%%%%%%%%%%%%%%%%%%%%%%%%%%%%%%%%%%%%%%%%%%%%%%%%%%%
%%%%%%%%%%%%%%%%%%%%%%%%%%%%%%%%%%%%%%%%%%%%%%%%%%%%%%%%%%%%%%%%%%%%%%%%%%%%%%%%
%%%%%%%%%%%%%%%%%%%%%%%%%%%%%%%%%%%%%%%%%%%%%%%%%%%%%%%%%%%%%%%%%%%%%%%%%%%%%%%%
\begin{mdframed}[style=darkQuesion]
17. Let $a,b$ be nonzero integers. Prove $\lrp{a,b}=1$ if and only if 
$\lrp{a+b,ab}=1$. 
\end{mdframed}

%%%%%%%%%%%%%%%%%%%%%%%%%%%%%%%%%%%%%%%%%%%%%%%%%%%%%%%%%%%%%%%%%%%%%%%%%%%%%%%%
\begin{mdframed}[style=darkAnswer,frametitle={Joe Starr}]
\begin{itemize}[align=left]
    \item [$\Rightarrow$]{
    We let $\lrp{a,b}=1$, then consider the $\lrp{a+b,ab}$. We assume 
    $\lrp{a+b,ab}=d$, with $d>1$. Since $d>1$ there must exist $p$ a prime such that
    $p\vert d$. This means that $p\vert a+b$ and $p \vert ab$. Consequently, either
    $p\vert a$ or $p\vert b$. WOLG we have $p\vert a$, and since $p\vert a+b$ it 
    must be that $p\vert b$. Finally, since $p\vert a $ and $p\vert b$, 
    $p\vert \lrp{a,b}$ a contradiction. So $\lrp{a+b,ab}=1$.
    }
    \item [$\Leftarrow$]{
    We let $\lrp{a+b,ab}=1$, then consider the $\lrp{a,b}$. We assume 
    $\lrp{a,b}=d$, with $d>1$. Since $d>1$ there must exist $p$ a prime such that
    $p\vert d$. This means that $p\vert a$ and $p\vert b$, further $p \vert ab$. 
    Since $p$ divides $a$ and $b$, we have $p\vert a+b$. Finally, since $p\vert ab$, 
    and $p\vert a+b$, $p\vert \lrp{a+b,ab}$, a contradiction so $\lrp{a,b}=1$.
    }
\end{itemize}
\end{mdframed}
\newpage
%%%%%%%%%%%%%%%%%%%%%%%%%%%%%%%%%%%%%%%%%%%%%%%%%%%%%%%%%%%%%%%%%%%%%%%%%%%%%%%%
%%%%%%%%%%%%%%%%%%%%%%%%%%%%%%%%%%%%%%%%%%%%%%%%%%%%%%%%%%%%%%%%%%%%%%%%%%%%%%%%
%%%%%%%%%%%%%%%%%%%%%%%%%%%%%%%%%%%%%%%%%%%%%%%%%%%%%%%%%%%%%%%%%%%%%%%%%%%%%%%%
%%%%%%%%%%%%%%%%%%%%%%%%%%%%%%%%%%%%%%%%%%%%%%%%%%%%%%%%%%%%%%%%%%%%%%%%%%%%%%%%
\begin{mdframed}[style=darkQuesion]
18. Let $a,b$ be nonzero integers with $\lrp{a,b}=1$. Compute $\lrp{a+b,a-b}$. 
\end{mdframed}

%%%%%%%%%%%%%%%%%%%%%%%%%%%%%%%%%%%%%%%%%%%%%%%%%%%%%%%%%%%%%%%%%%%%%%%%%%%%%%%%
\begin{mdframed}[style=darkAnswer,frametitle={Joe Starr}]
    
\end{mdframed}
\newpage
%%%%%%%%%%%%%%%%%%%%%%%%%%%%%%%%%%%%%%%%%%%%%%%%%%%%%%%%%%%%%%%%%%%%%%%%%%%%%%%%
%%%%%%%%%%%%%%%%%%%%%%%%%%%%%%%%%%%%%%%%%%%%%%%%%%%%%%%%%%%%%%%%%%%%%%%%%%%%%%%%
%%%%%%%%%%%%%%%%%%%%%%%%%%%%%%%%%%%%%%%%%%%%%%%%%%%%%%%%%%%%%%%%%%%%%%%%%%%%%%%%
%%%%%%%%%%%%%%%%%%%%%%%%%%%%%%%%%%%%%%%%%%%%%%%%%%%%%%%%%%%%%%%%%%%%%%%%%%%%%%%%
\begin{mdframed}[style=darkQuesion]
19. Let $a$ and $b$ be positive integers, and let $m$ be an integer such that 
$ab=m\lrp{a,b}$. Without using the prime factorization theorem, prove that 
$\lrp{a,b}\lrb{a,b}=ab$ by verifying that $m$ satisfies the necessary properties
of $\lrb{a,b}$.
\end{mdframed}

%%%%%%%%%%%%%%%%%%%%%%%%%%%%%%%%%%%%%%%%%%%%%%%%%%%%%%%%%%%%%%%%%%%%%%%%%%%%%%%%
\begin{mdframed}[style=darkAnswer,frametitle={Joe Starr}]
    
\end{mdframed}
\newpage
%%%%%%%%%%%%%%%%%%%%%%%%%%%%%%%%%%%%%%%%%%%%%%%%%%%%%%%%%%%%%%%%%%%%%%%%%%%%%%%%
%%%%%%%%%%%%%%%%%%%%%%%%%%%%%%%%%%%%%%%%%%%%%%%%%%%%%%%%%%%%%%%%%%%%%%%%%%%%%%%%
%%%%%%%%%%%%%%%%%%%%%%%%%%%%%%%%%%%%%%%%%%%%%%%%%%%%%%%%%%%%%%%%%%%%%%%%%%%%%%%%
%%%%%%%%%%%%%%%%%%%%%%%%%%%%%%%%%%%%%%%%%%%%%%%%%%%%%%%%%%%%%%%%%%%%%%%%%%%%%%%%
\begin{mdframed}[style=darkQuesion]
20. A positive integer $a$ is called a square if $a=n^2$ for some $n\in \Z$. 
Show that the integer $a>1$ is a square if and only if every exponent in its 
prime factorization is even. 
\end{mdframed}

%%%%%%%%%%%%%%%%%%%%%%%%%%%%%%%%%%%%%%%%%%%%%%%%%%%%%%%%%%%%%%%%%%%%%%%%%%%%%%%%
\begin{mdframed}[style=darkAnswer,frametitle={Joe Starr}]
Let $a\in \Z$ be a square. Since $a$ is a square by definiton there exists 
a $n$ such that $nn=a$. Now by the fundamental theorem of arythmic we know $n$
has a prime factorization,written $p_1^{n_1}\cdots p_k^{n_k}$. If we consider 
$nn$, we have $nn = \lrp{p_1^{n_1}\cdots p_k^{n_k}}
                    \lrp{p_1^{n_1}\cdots p_k^{n_k}}$
, by combining terms we can see that $nn= \lrp{p_1^{2n_1}\cdots p_k^{2n_k}}$, as
desired. 
\end{mdframed}
\newpage
%%%%%%%%%%%%%%%%%%%%%%%%%%%%%%%%%%%%%%%%%%%%%%%%%%%%%%%%%%%%%%%%%%%%%%%%%%%%%%%%
%%%%%%%%%%%%%%%%%%%%%%%%%%%%%%%%%%%%%%%%%%%%%%%%%%%%%%%%%%%%%%%%%%%%%%%%%%%%%%%%
%%%%%%%%%%%%%%%%%%%%%%%%%%%%%%%%%%%%%%%%%%%%%%%%%%%%%%%%%%%%%%%%%%%%%%%%%%%%%%%%
%%%%%%%%%%%%%%%%%%%%%%%%%%%%%%%%%%%%%%%%%%%%%%%%%%%%%%%%%%%%%%%%%%%%%%%%%%%%%%%%
\begin{mdframed}[style=darkQuesion]
23. Let $p$ and $q$ be prime numbers. Prove that $pq+1$ is a square if and only 
if $p$ and $q$ are twin primes.
\end{mdframed}

%%%%%%%%%%%%%%%%%%%%%%%%%%%%%%%%%%%%%%%%%%%%%%%%%%%%%%%%%%%%%%%%%%%%%%%%%%%%%%%%
\begin{mdframed}[style=darkAnswer,frametitle={Joe Starr}]
    
\end{mdframed}
\newpage
%%%%%%%%%%%%%%%%%%%%%%%%%%%%%%%%%%%%%%%%%%%%%%%%%%%%%%%%%%%%%%%%%%%%%%%%%%%%%%%%
%%%%%%%%%%%%%%%%%%%%%%%%%%%%%%%%%%%%%%%%%%%%%%%%%%%%%%%%%%%%%%%%%%%%%%%%%%%%%%%%
%%%%%%%%%%%%%%%%%%%%%%%%%%%%%%%%%%%%%%%%%%%%%%%%%%%%%%%%%%%%%%%%%%%%%%%%%%%%%%%%
%%%%%%%%%%%%%%%%%%%%%%%%%%%%%%%%%%%%%%%%%%%%%%%%%%%%%%%%%%%%%%%%%%%%%%%%%%%%%%%%
\begin{mdframed}[style=darkQuesion]
26. Prove that if $a>1$, then there is a prime $p$ with $a<p\leq a!+1$.
\end{mdframed}

%%%%%%%%%%%%%%%%%%%%%%%%%%%%%%%%%%%%%%%%%%%%%%%%%%%%%%%%%%%%%%%%%%%%%%%%%%%%%%%%
\begin{mdframed}[style=darkAnswer,frametitle={Joe Starr}]
    
\end{mdframed}
\newpage
%%%%%%%%%%%%%%%%%%%%%%%%%%%%%%%%%%%%%%%%%%%%%%%%%%%%%%%%%%%%%%%%%%%%%%%%%%%%%%%%
%%%%%%%%%%%%%%%%%%%%%%%%%%%%%%%%%%%%%%%%%%%%%%%%%%%%%%%%%%%%%%%%%%%%%%%%%%%%%%%%
%%%%%%%%%%%%%%%%%%%%%%%%%%%%%%%%%%%%%%%%%%%%%%%%%%%%%%%%%%%%%%%%%%%%%%%%%%%%%%%%
%%%%%%%%%%%%%%%%%%%%%%%%%%%%%%%%%%%%%%%%%%%%%%%%%%%%%%%%%%%%%%%%%%%%%%%%%%%%%%%%
\begin{mdframed}[style=darkQuesion]
29. Show that $\log{2}/\log{3}$ is not a rational number. 
\end{mdframed}

%%%%%%%%%%%%%%%%%%%%%%%%%%%%%%%%%%%%%%%%%%%%%%%%%%%%%%%%%%%%%%%%%%%%%%%%%%%%%%%%
\begin{mdframed}[style=darkAnswer,frametitle={Joe Starr}]
    
\end{mdframed}



\clearpage
%Compile Chapter 2

\section{Functions}  
\subsection{Functions}
%%%%%%%%%%%%%%%%%%%%%%%%%%%%%%%%%%%%%%%%%%%%%%%%%%%%%%%%%%%%%%%%%%%%%%%%%%%%%%%%
%%%%%%%%%%%%%%%%%%%%%%%%%%%%%%%%%%%%%%%%%%%%%%%%%%%%%%%%%%%%%%%%%%%%%%%%%%%%%%%%
%%%%%%%%%%%%%%%%%%%%%%%%%%%%%%%%%%%%%%%%%%%%%%%%%%%%%%%%%%%%%%%%%%%%%%%%%%%%%%%%
%%%%%%%%%%%%%%%%%%%%%%%%%%%%%%%%%%%%%%%%%%%%%%%%%%%%%%%%%%%%%%%%%%%%%%%%%%%%%%%%
\begin{mdframed}[style=darkQuesion]
  1.   In each of the following parts, determine whether the given function is 
  1:1 and whether it is onto. 
\begin{itemize}
  \item [(a)]{
    $f:\R\to \R; \fof{x}=x+3$
  }
  \item [(b)]{
    $f:\C\to \C; \fof{x}=x^2+2x+1$
  }
  \item [(c]{
    $f:\Z_n\to \Z_n; \fof{\lrb{x}_n}=\lrb{mx+b}_n, \text{where} m,b\in \Z$
  }
  \item [(d)]{
    $f:\R^+\to \R; \fof{x}=\ln x$
  }
\end{itemize}
\end{mdframed}

%%%%%%%%%%%%%%%%%%%%%%%%%%%%%%%%%%%%%%%%%%%%%%%%%%%%%%%%%%%%%%%%%%%%%%%%%%%%%%%%
\begin{mdframed}[style=darkAnswer,frametitle={Joe Starr}]
  \begin{itemize}
    \item [(a)]{ We can see that $\fof{x}=x+3$ then $\fofinv{x}=x-3$, 
    $\fof{\fofinv{x}}=\lrp{x-3}+3=x$. Showing $f$ is a bijection.
    }
    \item [(b)]{
      \begin{itemize}
        \item[1:1:]{\hspace{.5in}\newline
          Assume $\fof{x}=25=\fof{y}$, we can see that if $x=4$, $\fof{x}=25$, 
          and $y=-6$, $\fof{y}=25$, showing $f$ not injective. 
        }
        \item[onto:]{\hspace{.5in}\newline
          Let $y\in \C$ we must now show there exists a $x \in \C$ such that
          $\fof{x}=y$. Consider $x=\sqrt{y}-1$, we can then take:
          \begin{align*}
            \fof{x} &= x^2+2x+1\\
            &= \lrp{\sqrt{y}-1}^2+2\lrp{\sqrt{y}-1}+1\\
            &= \lrp{\sqrt{y}-1}^2+2\sqrt{y}-2+1\\
            &= 1 - 2 \sqrt{y} + y+2\sqrt{y}-1\\
            &=  y\\
          \end{align*}
          showing $f$ surjective.  
        }
      \end{itemize}
    }
    \item [(c)]{
    Consider $\fofinv{x}=\lrb{\lrp{y-b}m^{-1}}_n$, now taking 
    $\fof{\fofinv{x}}$
    \begin{align*}
      \fof{\fofinv{x}} &= \lrb{m\lrp{x-b}m^{-1}+b}_n\\
      &= \lrb{\lrp{x-b}+b}_n\\
      &= \lrb{x}_n\\
    \end{align*}
    showing $f$ a bijection.
    }
    \item [(d)]{
      $f:\R^+\to \R; \fof{x}=\ln x$
      If we take $\fofinv{x}=e^x$, $\fof{\fofinv{x}}=\ln e^x = x$, showing $f$
      a bijection. 
    }
  \end{itemize}
\end{mdframed}
  \newpage
%%%%%%%%%%%%%%%%%%%%%%%%%%%%%%%%%%%%%%%%%%%%%%%%%%%%%%%%%%%%%%%%%%%%%%%%%%%%%%%%
%%%%%%%%%%%%%%%%%%%%%%%%%%%%%%%%%%%%%%%%%%%%%%%%%%%%%%%%%%%%%%%%%%%%%%%%%%%%%%%%
%%%%%%%%%%%%%%%%%%%%%%%%%%%%%%%%%%%%%%%%%%%%%%%%%%%%%%%%%%%%%%%%%%%%%%%%%%%%%%%%
%%%%%%%%%%%%%%%%%%%%%%%%%%%%%%%%%%%%%%%%%%%%%%%%%%%%%%%%%%%%%%%%%%%%%%%%%%%%%%%%
\begin{mdframed}[style=darkQuesion]
  3.   For each 1:1 and onto function in Exercise 1, find the inverse of the 
  function
\begin{itemize}
  \item [(a)]{
    $f:\R\to \R; \fof{x}=x+3$
  }
  \item [(b)]{
    $f:\C\to \C; \fof{x}=x^2+2x+1$
  }
  \item [(c]{
    $f:\Z_n\to \Z_n; \fof{\lrb{x}_n}=\lrb{mx+b}_n, \text{where} m,b\in \Z$
  }
  \item [(d)]{
    $f:\R^+\to \R; \fof{x}=\ln x$
  }
\end{itemize}
\end{mdframed}

%%%%%%%%%%%%%%%%%%%%%%%%%%%%%%%%%%%%%%%%%%%%%%%%%%%%%%%%%%%%%%%%%%%%%%%%%%%%%%%%
\begin{mdframed}[style=darkAnswer,frametitle={Joe Starr}]
  \begin{multicols}{2}
    \begin{itemize}
    \item [(a)]{
        see question 1
    }
    \item [(b)]{
      not a bijection
    }
    \item [(c]{
      see question 1
      }
      \item [(d)]{
      see question 1
      }
    \end{itemize}
  \end{multicols}
  \end{mdframed}
\newpage
%%%%%%%%%%%%%%%%%%%%%%%%%%%%%%%%%%%%%%%%%%%%%%%%%%%%%%%%%%%%%%%%%%%%%%%%%%%%%%%%
%%%%%%%%%%%%%%%%%%%%%%%%%%%%%%%%%%%%%%%%%%%%%%%%%%%%%%%%%%%%%%%%%%%%%%%%%%%%%%%%
%%%%%%%%%%%%%%%%%%%%%%%%%%%%%%%%%%%%%%%%%%%%%%%%%%%%%%%%%%%%%%%%%%%%%%%%%%%%%%%%
%%%%%%%%%%%%%%%%%%%%%%%%%%%%%%%%%%%%%%%%%%%%%%%%%%%%%%%%%%%%%%%%%%%%%%%%%%%%%%%%
\begin{mdframed}[style=darkQuesion]
  4.   For each 1:1 and onto function in Exercise 2, find the inverse of the 
  function
\begin{itemize}
  \item [(a)]{
    $f:\R\to \R; \fof{x}=x^2$
  }
  \item [(b)]{
    $f:\C\to \C; \fof{x}=x^2$
  }
  \item [(c]{
    $f:\R^+\to \R^+; \fof{x}=x^2$
  }
  \item [(d)]{
    $f:\R^+\to \R^+; \fof{x}=
    \begin{cases}
      x &\text{if }x\text{ is rational}\\
      x^2 &\text{if }x\text{ is irrational}\\
    \end{cases}$
  }
\end{itemize}
\end{mdframed}

%%%%%%%%%%%%%%%%%%%%%%%%%%%%%%%%%%%%%%%%%%%%%%%%%%%%%%%%%%%%%%%%%%%%%%%%%%%%%%%%
\begin{mdframed}[style=darkAnswer,frametitle={Joe Starr}]
  \begin{itemize}
    \item [(a)]{
      Not a bijection
    }
    \item [(b)]{
      Not a bijection 
    }
    \item [(c]{
      $\fofinv{x}=+\sqrt{x}$
    }
    \item [(d)]{
      $f:\R^+\to \R^+; \fofinv{x}=
      \begin{cases}
        x &\text{if }x\text{ is rational}\\
        +\sqrt{x} &\text{if }x\text{ is irrational}\\
      \end{cases}$
    }
  \end{itemize}
\end{mdframed}
\newpage
%%%%%%%%%%%%%%%%%%%%%%%%%%%%%%%%%%%%%%%%%%%%%%%%%%%%%%%%%%%%%%%%%%%%%%%%%%%%%%%%
%%%%%%%%%%%%%%%%%%%%%%%%%%%%%%%%%%%%%%%%%%%%%%%%%%%%%%%%%%%%%%%%%%%%%%%%%%%%%%%%
%%%%%%%%%%%%%%%%%%%%%%%%%%%%%%%%%%%%%%%%%%%%%%%%%%%%%%%%%%%%%%%%%%%%%%%%%%%%%%%%
%%%%%%%%%%%%%%%%%%%%%%%%%%%%%%%%%%%%%%%%%%%%%%%%%%%%%%%%%%%%%%%%%%%%%%%%%%%%%%%%
\begin{mdframed}[style=darkQuesion]
  13. Let $f:A\to B$ be a function, and let $\fof{A}=\lrs{\fof{a}\vert a\in A}$
  be the image of $f$. Show that $f$is onto if and only if $\fof{A}=B$.
\end{mdframed}

%%%%%%%%%%%%%%%%%%%%%%%%%%%%%%%%%%%%%%%%%%%%%%%%%%%%%%%%%%%%%%%%%%%%%%%%%%%%%%%%
\begin{mdframed}[style=darkAnswer,frametitle={Joe Starr}]
Let $\fof{A}=B$, select $y\in B$, since $y\in B$ we have $y\in \fof{A}$, that 
means there exists $a\in A$ such that $\fof{a}=y$. Showing $f$ surjective. 
Let $\fof{A}\neq B$, let $y\in B$, such that $y\notin B \cap \fof{A}$. Since we
have $y\notin \fof{A}$ we have no $a\in A$ that maps to $y$ showing $f$ not 
surjective, as desired.  
\end{mdframed}
\newpage
%%%%%%%%%%%%%%%%%%%%%%%%%%%%%%%%%%%%%%%%%%%%%%%%%%%%%%%%%%%%%%%%%%%%%%%%%%%%%%%%
%%%%%%%%%%%%%%%%%%%%%%%%%%%%%%%%%%%%%%%%%%%%%%%%%%%%%%%%%%%%%%%%%%%%%%%%%%%%%%%%
%%%%%%%%%%%%%%%%%%%%%%%%%%%%%%%%%%%%%%%%%%%%%%%%%%%%%%%%%%%%%%%%%%%%%%%%%%%%%%%%
%%%%%%%%%%%%%%%%%%%%%%%%%%%%%%%%%%%%%%%%%%%%%%%%%%%%%%%%%%%%%%%%%%%%%%%%%%%%%%%%
\begin{mdframed}[style=darkQuesion]
  15. Let $f:A\to B$ and $g:B\to C$ be functions. Prove that if $g\circ f$ is 
  1:1, then $f$ is 1:1, and that if $g\circ f$ is onto $g$ is onto.
\end{mdframed}

%%%%%%%%%%%%%%%%%%%%%%%%%%%%%%%%%%%%%%%%%%%%%%%%%%%%%%%%%%%%%%%%%%%%%%%%%%%%%%%%
\begin{mdframed}[style=darkAnswer,frametitle={Joe Starr}]
Let $g\circ f$ be injective, but $f$ not injective. Since $f$ is not injective
$\exists a,x\in A$ such tha $a\neq x$ but $\fof{x}=\fof{a}$. We consider 
$g\circ f\lrp{a}$ and $g\circ f\lrp{x}$, since $\fof{x}=\fof{a}$ it must be that
$\gof{\fof{x}}=\gof{\fof{a}}$. This means that with $a\neq b$, 
$\gof{\fof{x}}=\gof{\fof{a}}$, making $g\circ f$ not injective a contradiction 
so $f$ injective. 

Let $g\circ f$ be surjective, but $g$ not surjective. Since $g$ not surjective
there exists some $c\in C$ such that $\gof{b}\neq c$ for all $b\in B$. However 
since $g\circ f$ surjective there exists $g\circ f\lrp{a}=c$ a contradiction, 
making $g$ surjective. 
\end{mdframed}
\newpage
%%%%%%%%%%%%%%%%%%%%%%%%%%%%%%%%%%%%%%%%%%%%%%%%%%%%%%%%%%%%%%%%%%%%%%%%%%%%%%%%
%%%%%%%%%%%%%%%%%%%%%%%%%%%%%%%%%%%%%%%%%%%%%%%%%%%%%%%%%%%%%%%%%%%%%%%%%%%%%%%%
%%%%%%%%%%%%%%%%%%%%%%%%%%%%%%%%%%%%%%%%%%%%%%%%%%%%%%%%%%%%%%%%%%%%%%%%%%%%%%%%
%%%%%%%%%%%%%%%%%%%%%%%%%%%%%%%%%%%%%%%%%%%%%%%%%%%%%%%%%%%%%%%%%%%%%%%%%%%%%%%%
\begin{mdframed}[style=darkQuesion]
  17. Let $f:A\to B$ be a function. Prove that $f$ is onto if and only if 
  $h\circ f = k \circ f$ implies $h=k$, for every set $C$ and all choices of 
  functions $h:B\to C$ and $k:B\to C$.
\end{mdframed}

%%%%%%%%%%%%%%%%%%%%%%%%%%%%%%%%%%%%%%%%%%%%%%%%%%%%%%%%%%%%%%%%%%%%%%%%%%%%%%%%
\begin{mdframed}[style=darkAnswer,frametitle={Joe Starr}]

\end{mdframed}
\newpage
%%%%%%%%%%%%%%%%%%%%%%%%%%%%%%%%%%%%%%%%%%%%%%%%%%%%%%%%%%%%%%%%%%%%%%%%%%%%%%%%
%%%%%%%%%%%%%%%%%%%%%%%%%%%%%%%%%%%%%%%%%%%%%%%%%%%%%%%%%%%%%%%%%%%%%%%%%%%%%%%%
%%%%%%%%%%%%%%%%%%%%%%%%%%%%%%%%%%%%%%%%%%%%%%%%%%%%%%%%%%%%%%%%%%%%%%%%%%%%%%%%
%%%%%%%%%%%%%%%%%%%%%%%%%%%%%%%%%%%%%%%%%%%%%%%%%%%%%%%%%%%%%%%%%%%%%%%%%%%%%%%%
\begin{mdframed}[style=darkQuesion]
  19. Let $f:A\to B$ be a function. Prove that $f$ is 1:1 if and only if 
  $f\circ h=f\circ k$ implies $h=k$, for every set $C$ and all choices of 
  functions $h:C\to A$ and $k:C\to A$.
\end{mdframed}
  
%%%%%%%%%%%%%%%%%%%%%%%%%%%%%%%%%%%%%%%%%%%%%%%%%%%%%%%%%%%%%%%%%%%%%%%%%%%%%%%%
\begin{mdframed}[style=darkAnswer,frametitle={Joe Starr}]

\end{mdframed}
\newpage
%%%%%%%%%%%%%%%%%%%%%%%%%%%%%%%%%%%%%%%%%%%%%%%%%%%%%%%%%%%%%%%%%%%%%%%%%%%%%%%%
%%%%%%%%%%%%%%%%%%%%%%%%%%%%%%%%%%%%%%%%%%%%%%%%%%%%%%%%%%%%%%%%%%%%%%%%%%%%%%%%
%%%%%%%%%%%%%%%%%%%%%%%%%%%%%%%%%%%%%%%%%%%%%%%%%%%%%%%%%%%%%%%%%%%%%%%%%%%%%%%%
%%%%%%%%%%%%%%%%%%%%%%%%%%%%%%%%%%%%%%%%%%%%%%%%%%%%%%%%%%%%%%%%%%%%%%%%%%%%%%%%
\subsection{Equivalence Relations}
%%%%%%%%%%%%%%%%%%%%%%%%%%%%%%%%%%%%%%%%%%%%%%%%%%%%%%%%%%%%%%%%%%%%%%%%%%%%%%%%
%%%%%%%%%%%%%%%%%%%%%%%%%%%%%%%%%%%%%%%%%%%%%%%%%%%%%%%%%%%%%%%%%%%%%%%%%%%%%%%%
%%%%%%%%%%%%%%%%%%%%%%%%%%%%%%%%%%%%%%%%%%%%%%%%%%%%%%%%%%%%%%%%%%%%%%%%%%%%%%%%
%%%%%%%%%%%%%%%%%%%%%%%%%%%%%%%%%%%%%%%%%%%%%%%%%%%%%%%%%%%%%%%%%%%%%%%%%%%%%%%%
\newpage
%%%%%%%%%%%%%%%%%%%%%%%%%%%%%%%%%%%%%%%%%%%%%%%%%%%%%%%%%%%%%%%%%%%%%%%%%%%%%%%%
%%%%%%%%%%%%%%%%%%%%%%%%%%%%%%%%%%%%%%%%%%%%%%%%%%%%%%%%%%%%%%%%%%%%%%%%%%%%%%%%
%%%%%%%%%%%%%%%%%%%%%%%%%%%%%%%%%%%%%%%%%%%%%%%%%%%%%%%%%%%%%%%%%%%%%%%%%%%%%%%%
%%%%%%%%%%%%%%%%%%%%%%%%%%%%%%%%%%%%%%%%%%%%%%%%%%%%%%%%%%%%%%%%%%%%%%%%%%%%%%%%
\subsection{Permutations} %1,2,3,4,5
%%%%%%%%%%%%%%%%%%%%%%%%%%%%%%%%%%%%%%%%%%%%%%%%%%%%%%%%%%%%%%%%%%%%%%%%%%%%%%%%
%%%%%%%%%%%%%%%%%%%%%%%%%%%%%%%%%%%%%%%%%%%%%%%%%%%%%%%%%%%%%%%%%%%%%%%%%%%%%%%%
%%%%%%%%%%%%%%%%%%%%%%%%%%%%%%%%%%%%%%%%%%%%%%%%%%%%%%%%%%%%%%%%%%%%%%%%%%%%%%%%
%%%%%%%%%%%%%%%%%%%%%%%%%%%%%%%%%%%%%%%%%%%%%%%%%%%%%%%%%%%%%%%%%%%%%%%%%%%%%%%%
\begin{mdframed}[style=darkQuesion]
  1. Consider the following Permutations in $S_7$.
  \begin{multicols}{2}
    $\sigma=  
    \begin{pmatrix}
      1 & 2 & 3 & 4 & 5 & 6 & 7\\
      3 & 2 & 5 & 4 & 6 & 1 & 7\\
    \end{pmatrix}$
    $\tau=  
    \begin{pmatrix}
      1 & 2 & 3 & 4 & 5 & 6 & 7\\
      2 & 1 & 5 & 7 & 4 & 6 & 3\\
    \end{pmatrix}$
  \end{multicols}
  \vspace{.25in}
  \begin{itemize}
    \item [(a)]{$\sigma\tau$
      
    }
    \item [(b)]{$\tau\sigma$

    }
    \item [(c)]{$\tau^2\sigma$

    }
    \item [(d)]{$\sigma^{-1}$

    }
    \item [(e)]{$\sigma\tau\sigma^{-1}$

    }
    \item [(f)]{$\tau^{-1}\sigma\tau$

    }
  \end{itemize}
\end{mdframed}
  
%%%%%%%%%%%%%%%%%%%%%%%%%%%%%%%%%%%%%%%%%%%%%%%%%%%%%%%%%%%%%%%%%%%%%%%%%%%%%%%%
\begin{mdframed}[style=darkAnswer,frametitle={Joe Starr}]
  \begin{itemize}
    \item [(a)]{$\sigma\tau$
      
    }
    \item [(b)]{$\tau\sigma$

    }
    \item [(c)]{$\tau^2\sigma$

    }
    \item [(d)]{$\sigma^{-1}$

    }
    \item [(e)]{$\sigma\tau\sigma^{-1}$

    }
    \item [(f)]{$\tau^{-1}\sigma\tau$

    }
  \end{itemize}
\end{mdframed}
\newpage
%%%%%%%%%%%%%%%%%%%%%%%%%%%%%%%%%%%%%%%%%%%%%%%%%%%%%%%%%%%%%%%%%%%%%%%%%%%%%%%%
%%%%%%%%%%%%%%%%%%%%%%%%%%%%%%%%%%%%%%%%%%%%%%%%%%%%%%%%%%%%%%%%%%%%%%%%%%%%%%%%
%%%%%%%%%%%%%%%%%%%%%%%%%%%%%%%%%%%%%%%%%%%%%%%%%%%%%%%%%%%%%%%%%%%%%%%%%%%%%%%%
%%%%%%%%%%%%%%%%%%%%%%%%%%%%%%%%%%%%%%%%%%%%%%%%%%%%%%%%%%%%%%%%%%%%%%%%%%%%%%%%
\begin{mdframed}[style=darkQuesion]
  2. Write each of the permutations $\sigma\tau, \tau\sigma, \tau^2\sigma, 
  \sigma^{-1}, \sigma\tau\sigma^{-1}, \text{ and } \tau^{-1}\sigma\tau$ in 
  Exercise 1 as a product of disjoint cycles. Write $\sigma$ anf $\tau$ as 
  products of transpositions. 
\end{mdframed}
  
%%%%%%%%%%%%%%%%%%%%%%%%%%%%%%%%%%%%%%%%%%%%%%%%%%%%%%%%%%%%%%%%%%%%%%%%%%%%%%%%
\begin{mdframed}[style=darkAnswer,frametitle={Joe Starr}]
  \begin{itemize}
    \item [(a)]{$\sigma\tau$
      
    }
    \item [(b)]{$\tau\sigma$

    }
    \item [(c)]{$\tau^2\sigma$

    }
    \item [(d)]{$\sigma^{-1}$

    }
    \item [(e)]{$\sigma\tau\sigma^{-1}$

    }
    \item [(f)]{$\tau^{-1}\sigma\tau$

    }
  \end{itemize}
\end{mdframed}
\newpage
%%%%%%%%%%%%%%%%%%%%%%%%%%%%%%%%%%%%%%%%%%%%%%%%%%%%%%%%%%%%%%%%%%%%%%%%%%%%%%%%
%%%%%%%%%%%%%%%%%%%%%%%%%%%%%%%%%%%%%%%%%%%%%%%%%%%%%%%%%%%%%%%%%%%%%%%%%%%%%%%%
%%%%%%%%%%%%%%%%%%%%%%%%%%%%%%%%%%%%%%%%%%%%%%%%%%%%%%%%%%%%%%%%%%%%%%%%%%%%%%%%
%%%%%%%%%%%%%%%%%%%%%%%%%%%%%%%%%%%%%%%%%%%%%%%%%%%%%%%%%%%%%%%%%%%%%%%%%%%%%%%%
\begin{mdframed}[style=darkQuesion]
  3. Write 
    $\begin{pmatrix}
      1 & 2 & 3 & 4 & 5 & 6 & 7 & 8 & 9 & 10\\
      3 & 4 & 10 & 5 & 7 & 8 & 2 & 6 & 9 & 1\\
    \end{pmatrix}$ as the product of disjoint cycles and as a product of 
    transpositions. Construct its associated diagram, find its inverse, and find 
    it's order. 
\end{mdframed}
  
%%%%%%%%%%%%%%%%%%%%%%%%%%%%%%%%%%%%%%%%%%%%%%%%%%%%%%%%%%%%%%%%%%%%%%%%%%%%%%%%
\begin{mdframed}[style=darkAnswer,frametitle={Joe Starr}]
\end{mdframed}
\newpage
%%%%%%%%%%%%%%%%%%%%%%%%%%%%%%%%%%%%%%%%%%%%%%%%%%%%%%%%%%%%%%%%%%%%%%%%%%%%%%%%
%%%%%%%%%%%%%%%%%%%%%%%%%%%%%%%%%%%%%%%%%%%%%%%%%%%%%%%%%%%%%%%%%%%%%%%%%%%%%%%%
%%%%%%%%%%%%%%%%%%%%%%%%%%%%%%%%%%%%%%%%%%%%%%%%%%%%%%%%%%%%%%%%%%%%%%%%%%%%%%%%
%%%%%%%%%%%%%%%%%%%%%%%%%%%%%%%%%%%%%%%%%%%%%%%%%%%%%%%%%%%%%%%%%%%%%%%%%%%%%%%%
\begin{mdframed}[style=darkQuesion]
  3. Find the oder of each of the following permutations. 
  \begin{itemize}
    \item [(a)] {
      $\begin{pmatrix}
        1 & 2 & 3 & 4 & 5 & 6\\
        6 & 4 & 5 & 3 & 2 & 1\\
      \end{pmatrix}$
    }
    \item [(b)] {
      $\begin{pmatrix}
        1 & 2 & 3 & 4 & 5 & 6 & 7 & 8\\
        4 & 6 & 7 & 5 & 1 & 8 & 2 & 3\\
      \end{pmatrix}$
    }
    \item [(c)] {
      $\begin{pmatrix}
        1 & 2 & 3 & 4 & 5 & 6 & 7 & 8 & 9\\
        5 & 9 & 8 & 7 & 3 & 4 & 6 & 1 & 2\\
      \end{pmatrix}$
    }
    \item [(d)] {
      $\begin{pmatrix}
        1 & 2 & 3 & 4 & 5 & 6 & 7 & 8 & 9 \\
        8 & 4 & 9 & 6 & 5 & 2 & 3 & 1 & 7\\
      \end{pmatrix}$
    }
  \end{itemize} 
\end{mdframed}
  
%%%%%%%%%%%%%%%%%%%%%%%%%%%%%%%%%%%%%%%%%%%%%%%%%%%%%%%%%%%%%%%%%%%%%%%%%%%%%%%%
\begin{mdframed}[style=darkAnswer,frametitle={Joe Starr}]
  \begin{itemize}
    \item [(a)] {
      $\begin{pmatrix}
        1 & 2 & 3 & 4 & 5 & 6\\
        6 & 4 & 5 & 3 & 2 & 1\\
      \end{pmatrix}$
    }
    \item [(b)] {
      $\begin{pmatrix}
        1 & 2 & 3 & 4 & 5 & 6 & 7 & 8\\
        4 & 6 & 7 & 5 & 1 & 8 & 2 & 3\\
      \end{pmatrix}$
    }
    \item [(c)] {
      $\begin{pmatrix}
        1 & 2 & 3 & 4 & 5 & 6 & 7 & 8 & 9\\
        5 & 9 & 8 & 7 & 3 & 4 & 6 & 1 & 2\\
      \end{pmatrix}$
    }
    \item [(d)] {
      $\begin{pmatrix}
        1 & 2 & 3 & 4 & 5 & 6 & 7 & 8 & 9 \\
        8 & 4 & 9 & 6 & 5 & 2 & 3 & 1 & 7\\
      \end{pmatrix}$
    }
  \end{itemize} 
\end{mdframed}
\newpage
%%%%%%%%%%%%%%%%%%%%%%%%%%%%%%%%%%%%%%%%%%%%%%%%%%%%%%%%%%%%%%%%%%%%%%%%%%%%%%%%
%%%%%%%%%%%%%%%%%%%%%%%%%%%%%%%%%%%%%%%%%%%%%%%%%%%%%%%%%%%%%%%%%%%%%%%%%%%%%%%%
%%%%%%%%%%%%%%%%%%%%%%%%%%%%%%%%%%%%%%%%%%%%%%%%%%%%%%%%%%%%%%%%%%%%%%%%%%%%%%%%
%%%%%%%%%%%%%%%%%%%%%%%%%%%%%%%%%%%%%%%%%%%%%%%%%%%%%%%%%%%%%%%%%%%%%%%%%%%%%%%%
\begin{mdframed}[style=darkQuesion]
  5. Let $3\leq m\leq n$. Calculate $\sigma\tau^{-1}$ for cycles 
  $\sigma= \lrp{1,2,\dots , m-1}$ and \\ $\tau= \lrp{1,2,\dots,m-1,m}$ in $S_n$.
\end{mdframed}
  
%%%%%%%%%%%%%%%%%%%%%%%%%%%%%%%%%%%%%%%%%%%%%%%%%%%%%%%%%%%%%%%%%%%%%%%%%%%%%%%%
\begin{mdframed}[style=darkAnswer,frametitle={Joe Starr}]
\end{mdframed}
\newpage

\clearpage
%Compile Chapter 3
\section{Section}  

I <3 my Wayne State Libraries! Do you? 

\subsection{A Subsection}

\clearpage
%Compile Chapter 4
\section{Section}  

I <3 my Wayne State Libraries! Do you? 

\subsection{A Subsection}

\clearpage
%Compile Chapter 5
\section{Section}  

I <3 my Wayne State Libraries! Do you? 

\subsection{A Subsection}

\clearpage
%Compile Chapter 6
\section{Section}  

I <3 my Wayne State Libraries! Do you? 

\subsection{A Subsection}

\clearpage
%Compile Chapter 7
\section{Section}  

I <3 my Wayne State Libraries! Do you? 

\subsection{A Subsection}

\clearpage
%Compile Chapter 8
\section{Section}  

I <3 my Wayne State Libraries! Do you? 

\subsection{A Subsection}

\clearpage
%Compile Chapter 9
\section{Section}  

I <3 my Wayne State Libraries! Do you? 

\subsection{A Subsection}

\clearpage
%Compile Chapter 10
\section{Section}  

I <3 my Wayne State Libraries! Do you? 

\subsection{A Subsection}

\clearpage

\end{document}